\subsubsection{Data Augmentation Results}
\label{sec:augmentation}

% Owner: Member 4. Bullet-level skeleton; expand to ~1--1.5 pages
% (this section can be the most visual one).

\paragraph{Experiment focus.}
We evaluate the augmentation families defined in \cref{sec:method:aug}
under the shared protocol of \cref{sec:exp:training,sec:exp:eval}.
Only the data pipeline changes; the model, optimizer, schedule, and
checkpoint-selection rule remain fixed.

\paragraph{Results.}
\draftnote{Expand this subsection into prose using the checklist below.}
\begin{itemize}
  \item Table: per augmentation (or per combination), best val acc /
        test acc / train--val gap / train time.
  \item Curves figure overlaying val accuracy per augmentation.
  \item Optional: a robustness probe (e.g. corruption noise) if time
        permits; otherwise skip.
\end{itemize}

\paragraph{Discussion.}
\label{sec:aug:discussion}
\draftnote{Answer the three Member 4 questions in short paragraphs.}
\begin{enumerate}
  \item Which augmentation is most effective?
  \item How big is the gap between simple and strong augmentation?
  \item Is data-level regularization more effective than
        parameter-level (cross-reference \cref{sec:wd,sec:dropout})?
\end{enumerate}
